% Generated draft: requires mathematical and source review.
\documentclass[11pt]{article}
\usepackage[margin=25mm]{geometry}
\usepackage{fontspec}
\setmainfont{Latin Modern Roman}
\usepackage{amsmath,amssymb,mathtools,unicode-math}
\setmathfont{Latin Modern Math}
\usepackage{enumitem,needspace,xurl,hyperref}
\setlength{\parindent}{0pt}
\setlength{\parskip}{5pt}
\setlength{\emergencystretch}{3em}
\newcommand{\N}{\mathbb{N}}
\newcommand{\Z}{\mathbb{Z}}
\newcommand{\Q}{\mathbb{Q}}
\newcommand{\R}{\mathbb{R}}
\newcommand{\C}{\mathbb{C}}
\newcommand{\F}{\mathbb{F}}
\newcommand{\A}{\mathbb{A}}
\newcommand{\PP}{\mathbb{P}}
\newcommand{\E}{\mathbb{E}}
\newcommand{\eps}{\varepsilon}
\newcommand{\dd}{\,\mathrm{d}}
\begin{document}
\section*{UnsolvedMath ID: 30005007}
\textbf{Source problem code:} catalogue entry
\subsection{Definitions and mathematical statement}
For terminology only, adopt the usual real polynomial interpolation convention; this convention does not recover the missing hypotheses of the excerpt. For an integer \(d\geq 1\), let
\[
\Pi_2^d=\{p\in\R[x_1,\ldots,x_d]:\deg p\leq 2\}.
\]
A finite set \(X\subset\R^d\) is \(2\)-correct if the evaluation map
\[
\operatorname{ev}_X:\Pi_2^d\longrightarrow\R^X,
\qquad p\longmapsto(p(x))_{x\in X},
\]
is bijective. In particular, \(|X|=\binom{d+2}{2}\). For each \(a\in X\), its fundamental Lagrange polynomial \(L_a\in\Pi_2^d\) satisfies \(L_a(a)=1\) and \(L_a(x)=0\) for every \(x\in X\setminus\{a\}\). Under this convention, \(X\) is a \(GC_2\) set if it is \(2\)-correct and every \(L_a\) is a product of two nonconstant affine linear polynomials. A maximal hyperplane means an affine hyperplane containing \(\binom{d+1}{2}\) nodes of \(X\).

\textbf{Recoverable content and unresolved statement.} The supplied record does not determine a standalone conjecture or problem. Its opening word refers to an omitted conjecture, and its final construction is unfinished. The following are descriptions of claims in the supplied excerpt, not independently verified assertions or a newly formulated conjecture.

For \(d=3\), the excerpt describes a \(GC_2\) configuration \(X\) with four maximal hyperplanes. It chooses a triple \(A_i=\{a_i,a_j,a_k\}\) and replaces an edge point \(b_{jk}\) by \(b'_{jk}\), producing \(X'\). The permitted location of \(b'_{jk}\) is corrupted in the supplied text and cannot be specified reliably. The excerpt claims that \(h_i\) loses maximality, the other \(h_j\) remain maximal, and \(X'\) remains \(GC_2\), with three maximal hyperplanes. It invokes an incompletely rendered invertibility condition, maximal-hyperplane factors for the fundamental polynomials at nodes other than \(a_i\), and an unchanged fundamental polynomial at \(a_i\) because the moved point passes from \(h_i\) into \(k_i\).

The excerpt then describes a second move associated with \(A_k=\{a_k,a_i,a_o\}\). It reports that the resulting \(X''\) is still \(2\)-correct but is not \(GC_2\): the fundamental polynomial at the common edge point \(b_{ik}\) no longer factors into linear factors. Its stated explanation concerns a point formerly in \(h_i\cup h_k\). The configurations, hyperplanes, and move conditions needed to express this claim precisely are absent.

The excerpt also mentions an abandoned attempt to obtain a counterexample in \(\R^4\), following a question from Jes\'us Carnicer, and attributes a counterexample for \(d=6\) to Apozyan, reference [1]. The latter construction begins by choosing a three-element set \(A_i=\{a_i,a_j,a_k\}\) for every \(i\in\{0,\ldots,d\}\), but its required conditions are cut off. The conjecture being contradicted and the full cited results are missing. Thus no complete mathematical target, with its hypotheses and quantifiers, can be recovered from this record.
\subsection{Short English statement}
This record preserves an incomplete discussion of higher-dimensional \(GC_2\) interpolation configurations. It reports a successful first point move, a failed second move, an abandoned attempt in \(\R^4\), and a counterexample attributed to Apozyan in dimension six. The omitted conjecture, missing configuration definitions, corrupted move condition, and truncated construction prevent a precise standalone problem statement. The supplied label ``open'' is a curation label, not a verified mathematical conclusion.
\subsection{Sources}
\begin{itemize}
\item Supplied local registry record: permanent UnsolvedMath ID \(30005007\), title \emph{Higher-Dimensional GC Set Counterexamples}, source problem code \texttt{catalogue entry}, in the Oberwolfach Reports Open Problems collection. The supplied proposed year is 2022; no proposer is identified.
\item Supplied source citation: \emph{Mini-Workshop: Interpolation, Approximation, and Algebra (2023)}, Oberwolfach Reports, pp.\ 383--428, DOI \texttt{10.4171/owr/2022/7}, \url{https://doi.org/10.4171/owr/2022/7}. The extracted passage includes a page marker for p.\ 395 and a caption for Figure 1. The figure itself is not supplied. The citation's parenthetical year 2023 and the supplied proposed year and DOI year 2022 are retained without reconciliation.
\item Supplied literature source: ``OWR interpolation report,'' \url{https://doi.org/10.4171/owr/2022/7}. Literature assessment checked at \texttt{2026-08-22}: this is supplied dated curation, not independently verified. It labels the record open because the incomplete passage does not state a complete proposition to which a definite truth value can be assigned. The supplied background expressly calls for independent verification of time-sensitive status and cited claims.
\item The supplied statement-verification field describes the passage as an account of known higher-dimensional counterexamples and unsuccessful approaches, rather than a new standalone open problem. It assigns statement status \texttt{unrecoverable}, promotes no clean statement, and retains the original record for provenance. This verification description is supplied metadata, not an independent verification performed here.
\item Internal reference [3] is credited with the dimension-three counterexample; internal reference [1] is attributed to Apozyan for the dimension-six counterexample. Their bibliographic entries are absent. The excerpt credits a question from Jes\'us Carnicer with exposing the failed second move; the speaker or author is not identified. Its suggested causal connection to Apozyan's construction is explicitly speculative.
\end{itemize}
\end{document}
